% =============================================================================
% GLOSSARY ENTRIES
% =============================================================================

% --- ACRONYMS ----------------------------------------------------------------
% Usage in text: \gls{qkd}  first use: "Quantum Key Distribution (QKD)"
% subsequent: "QKD"
% \glspl{qkd}  plural form

\newacronym{qkd}{QKD}{Quantum Key Distribution}
\newacronym{bb84}{BB84}{Bennett--Brassard 1984 Protocol}
\newacronym{b92}{B92}{Bennett 1992 Protocol}
\newacronym{e91}{E91}{Ekert 1991 Protocol}
\newacronym{qber}{QBER}{Quantum Bit Error Rate}
\newacronym{pqc}{PQC}{Post-Quantum Cryptography}
\newacronym{otp}{OTP}{One-Time Pad}
\newacronym{prng}{PRNG}{Pseudo-Random Number Generator}
\newacronym{epe}{EPE}{Entanglement Purification via Eavesdropping}
\newacronym{des}{DES}{Data Encryption Standard}
\newacronym{aes}{AES}{Advanced Encryption Standard}
\newacronym{chsh}{CHSH}{Clauser--Horne--Shimony--Holt}
\newacronym{edi}{EDI}{Equality, Diversity, and Inclusion}
\newacronym{epr}{EPR}{Einstein--Podolsky--Rosen}
\newacronym{cv}{CV}{Continuous-Variable}
\newacronym{dv}{DV}{Discrete-Variable}
\newacronym[sort=pandm]{pandm}{P\&M}{Prepare-and-Measure}
\newacronym{eb}{EB}{Entanglement-Based}
\newacronym{mdi}{MDI}{Measurement-Device-Independent}
\newacronym{cc}{CC}{Classical Channel}
\newacronym{oop}{OOP}{Object-Oriented Programming}
\newacronym{bsc}{BSC}{Binary Symmetric Channel}
\newacronym{sdk}{SDK}{Software Development Kit}
\newacronym{ide}{IDE}{Integrated Development Environment}
\newacronym{qchan}{QC}{Quantum Channel}
\newacronym{qcomp}{QC}{Quantum Computer}
\newacronym{skr}{SKR}{Secure Key Rate}
\newacronym{p2p}{P2P}{Peer-to-Peer}

% --- SYMBOLS (Nomenclature) --------------------------------------------------
% Usage in text: \gls{sym:qber}  renders the symbol with a hyperlink
%                \glsdesc{sym:qber} renders just the description

\newglossaryentry{sym:qber}{
    type    = symbols,
    name    = {\ensuremath{\mathrm{QBER}}},
    description = {Quantum Bit Error Rate --- ratio of erroneous bits to total
                   bits received over the quantum channel},
    sort    = {QBER}
}

\newglossaryentry{sym:rsec}{
    type    = symbols,
    name    = {\ensuremath{R_{\mathrm{sec}}}},
    description = {Secret-key generation rate (bits per channel use)},
    sort    = {Rsec}
}

\newglossaryentry{sym:theta}{
    type    = symbols,
    name    = {\ensuremath{\theta}},
    description = {Polarisation angle of a photon (radians)},
    sort    = {theta}
}

\newglossaryentry{sym:epsilon}{
    type    = symbols,
    name    = {\ensuremath{\varepsilon}},
    description = {Security parameter --- maximum acceptable eavesdropper
                   information leakage},
    sort    = {epsilon}
}

% --- GENERAL GLOSSARY TERMS --------------------------------------------------
% Usage in text: \gls{entanglement}

\newglossaryentry{entanglement}{
    name        = {entanglement},
    description = {A quantum mechanical phenomenon in which two or more particles
                   share a correlated quantum state, such that measuring one
                   particle instantaneously determines properties of the other
                   regardless of the distance separating them}
}

\newglossaryentry{qubit}{
    name        = {qubit},
    description = {The fundamental unit of quantum information; a two-level
                   quantum system that can exist in a superposition of states
                   \ensuremath{\ket{0}} and \ensuremath{\ket{1}}}
}

\newglossaryentry{sifting}{
    name        = {sifting},
    description = {The post-transmission step in QKD in which Alice and Bob
                   publicly compare measurement bases and discard bits where
                   incompatible bases were used}
}

\newglossaryentry{reconciliation}{
    name        = {reconciliation},
    description = {An error-correction stage in QKD that resolves discrepancies
                   between Alice's and Bob's sifted keys over an authenticated
                   classical channel}
}

\newglossaryentry{amplification}{
    name        = {privacy amplification},
    description = {A post-processing step that compresses the reconciled key to
                   reduce any partial information an eavesdropper may have
                   obtained, producing a shorter but provably secure final key}
}


% =============================================================================
% ENTRIES ADDED FOR: Introduction paragraph
% =============================================================================

% --- New Acronyms ------------------------------------------------------------

\newacronym{rsa}{RSA}{Rivest--Shamir--Adleman}
% First use in text: "Rivest--Shamir--Adleman (RSA)"
% Subsequent uses:   "RSA"

\newacronym{ecc}{ECC}{Elliptic Curve Cryptography}
% First use in text: "Elliptic Curve Cryptography (ECC)"
% Subsequent uses:   "ECC"

% --- New General Terms -------------------------------------------------------

\newglossaryentry{intfact}{
    name        = {integer factorisation},
    description = {The computational problem of decomposing a composite integer
                   into its prime factors. No classical algorithm is known to
                   solve this in polynomial time for large integers, forming the
                   security basis of RSA}
}

\newglossaryentry{disclog}{
    name        = {discrete logarithm},
    description = {The problem of finding the exponent \ensuremath{x} such that
                   \ensuremath{g^x \equiv h \pmod{p}} for given \ensuremath{g},
                   \ensuremath{h}, and prime \ensuremath{p}. Believed to be
                   classically hard; underlies the security of ECC and
                   Diffie--Hellman key exchange}
}

\newglossaryentry{shors}{
    name        = {Shor's algorithm},
    description = {A quantum algorithm, published by Peter Shor in 1994, that
                   solves integer factorisation and the discrete logarithm in
                   polynomial time on a fault-tolerant quantum computer,
                   breaking RSA and ECC}
}

\newglossaryentry{itsec}{
    name        = {information-theoretic security},
    description = {A security guarantee that holds against an adversary with
                   unlimited computational power, relying on mathematical
                   impossibility rather than assumed hardness of a problem.
                   QKD achieves this via the laws of quantum mechanics}
}

\newglossaryentry{qea}{
    name        = {Quantum Economic Advantage},
    description = {The threshold at which the practical cost of executing a
                   quantum attack --- accounting for hardware, error correction,
                   and runtime --- becomes economically feasible for a
                   sufficiently motivated adversary}
}

